\section{Discussion, extensions, and limitations}

The efficiency analysis links three aspects of private trial design: optimization for the treatment-effect contrast, adaptation to unknown arm means, and the number of outputs required by an attaining stationary release. Under the interior binary-trial model, known assignment probability, and fixed positive privacy level, \cref{obj:thm:finite-oracle} characterizes the stationary optimum through a finite staircase program. After fixing any measurable strong-saddle selector satisfying the conditions in \cref{obj:def:pilot-estimator}, \cref{obj:thm:sequential-minimax} identifies the optimized variance as the exact regular local asymptotic squared-error constant for sequential private procedures. Its solution determines the attainable first-order precision for estimating the treatment effect.

Contrast-specific optimization allocates information according to the scalar estimand while accounting for the unknown arm means through nuisance profiling. The balanced case in \cref{obj:thm:balanced-reduction} makes this principle especially transparent. With balanced assignment and complementary arm means, binary randomized response applied to the treatment--outcome sign attains the optimized contrast variance. More generally, \cref{obj:thm:five-output-upper-bound} guarantees an attaining staircase release with at most five active outputs at every admissible interior design point.

The private pilot makes parameter-dependent optimization operational. Under \cref{obj:ass:pilot-diverges,obj:ass:pilot-sublinear}, its size \(m_n\) grows while its fraction \(m_n/n\) vanishes. The increasing pilot learns the arm means through randomized-response releases, and the remaining subjects use the channel selected at the pilot estimate. Conditional on the pilot history, that channel and its affine correction are fixed. The exact conditional centering in \cref{obj:thm:pilot-attainment} allows the pilot to guide release selection while the main sample determines the first-order variance. With the stated strong-saddle selection conditions, every pilot schedule satisfying these restrictions yields the optimized normal limit and consistent studentization.

Support changes describe how the attaining stationary experiment varies with outcome prevalence. At assignment probability \(p=3/10\) and privacy level \(\varepsilon=\log 3\), \cref{obj:thm:support-transition} establishes open neighborhoods in which the minimum attaining cardinality is respectively five and three. These conclusions concern stationary channels at specified parameter points; the pilot construction adapts by selecting a channel from private observations. Its inference guarantee applies pointwise throughout the interior domain, including points at which an optimal support changes. The measurable strong-saddle conditions in \cref{obj:thm:pilot-attainment} govern selection at these points, and the empirical variance consistently estimates the optimized variance appearing in the limiting distribution. Thus support changes are compatible with pointwise asymptotic Wald coverage.

For prospective trial design, these results connect endpoint prevalence, privacy, and recruitment requirements through the optimized variance. The attendance-threshold endpoint provides one prospective binary outcome to which this analysis could apply, with the observation window and recording rules fixed before assignment. Planning arm means determine precision scenarios, while the private pilot supplies the estimates used during the trial. The comparison in \cref{obj:thm:laplace-comparison} translates the variance difference between the efficient procedure and the specified custom-Laplace sample mean into asymptotic interval-length and sample-size quantities.

\subsection{Limitations and future work}

A complete classification of optimal supports remains open. The five-output upper bound, the certified three- and five-output neighborhoods, and the balanced binary reduction establish specific features of the efficiency surface. Further work could characterize all attaining supports, ties between optimizers, and the boundaries separating support regions. The pointwise coverage result also leaves uniform coverage over parameter sequences approaching such boundaries as a separate question.

The attainment result applies to any measurable strong-saddle selector satisfying the stated conditions. Constructing a numerical reference solver with a deterministic tie-breaking rule and an independently checkable numerical certificate for a computed selector is a separate computational task.

The asymptotic analysis assumes interior arm means, a fixed assignment probability strictly between zero and one, and a fixed positive privacy level. Boundary arm means require an analysis of the corresponding local experiments. Assignment probabilities or privacy budgets that change with sample size may also alter pilot informativeness, the appropriate normalization, and the attainable variance. These regimes require conditions tailored to their rates of change.

Broader outcome models would require new record representations and information calculations, together with an appropriate private learning procedure for release selection. Repeated-user interaction introduces an additional distinction: the sequential protocols here adapt across independently sampled subjects, each supplying a private record. Multiple releases from the same individual require a privacy accounting and statistical model that accommodate dependence within that individual's transcript.

Finally, \cref{obj:thm:laplace-comparison} compares the efficient procedure with the sample mean of the specified custom-Laplace releases. Comparing that sample mean with full-likelihood estimation under the same release would clarify the contribution of estimator choice within the comparator experiment. Such an analysis would complement the present comparison of release-and-estimator procedures.
